\documentclass[12pt]{article}
\usepackage[T1]{fontenc}
\usepackage[utf8]{inputenc}
\usepackage{lmodern}
\usepackage{textcomp}
\usepackage{enumitem}
\usepackage{geometry}
\usepackage{graphicx}
\usepackage{amsmath, amssymb}
\geometry{margin=1in}
\begin{document}
\begin{center}
History - Undergraduate Level Assessment 12 \
Total Marks: 40 \
Answer all questions.
\end{center}

\section*{Multiple Choice (10 marks)}
\begin{enumerate}[label=\arabic*.]
\item By 1900 B.P., the central elements of Mayan culture were in place. These did NOT include:
\begin{enumerate}[label=(\alph*)]
    \item a hieroglyphic written language and calendar.
    \item stratified societies ruled by kings.
    \item a sophisticated bronze producing industry.
    \item construction of large-scale pyramids.
    \item usage of gold currency.
\end{enumerate}
\item Virtually all of the crops important in the European Neolithic:
\begin{enumerate}[label=(\alph*)]
    \item were primarily animal products, not plant-based.
    \item were imported from other regions.
    \item were domesticated before groups became sedentary.
    \item were domesticated at the same time at communal farms.
    \item were mainly wild crops that were not domesticated until the Bronze Age.
\end{enumerate}
\item The native peoples of the northwest coast of North America were:
\begin{enumerate}[label=(\alph*)]
    \item simple foragers.
    \item maize agriculturalists.
    \item conquered by the Aztec.
    \item dependent on hunting large game.
    \item primarily seafaring traders.
\end{enumerate}
\item The Hopewell were complex rank societies, but they were not a state. Hopewell culture lacked all of the following elements EXCEPT:
\begin{enumerate}[label=(\alph*)]
    \item formal government
    \item monumental earthworks
    \item urban centers
    \item a permanent military
\end{enumerate}
\item Archaeological evidence for the collapse of civilizations suggests which of the following is the most important variable?
\begin{enumerate}[label=(\alph*)]
    \item the existence of a written language
    \item the number of years the civilization has existed
    \item the civilization's religious beliefs and practices
    \item whether warfare can be ended
    \item the size of the civilization's population
\end{enumerate}
\end{enumerate}

\section*{True / False (10 marks)}
\textit{Answer True or False}
\begin{enumerate}[label=\arabic*.]
\setcounter{enumi}{5}
\item True or False: To secure further loans, Westinghouse was forced to revisit Tesla's AC patent, which bankers considered a financial strain on the company (at that point Westinghouse had paid out a different answer in licenses and royalties to Tesla, Brown, and Peck).
\item True or False: The above "history of economics" reflects modern economic textbooks and this means that the last stage of a science is represented as the culmination of its history (Kuhn, 1962). The "invisible hand" mentioned in a lost page in the middle of a chapter in the middle of the to "Wealth of Nations", 1776, advances as Smith's central message.
\item True or False: Some Presbyterian traditions adopt only the Westminster Confession of Faith as the doctrinal standard to which teaching elders are required to subscribe, in contrast to the Larger and Shorter catechisms, which are approved for use in instruction.
\item True or False: Elsewhere, remnants of the medieval water supply system devised by the friars can still be seen today. Constructed in 1290, the system carried water from Conduit Head (remnants of which survive near Hill Lane, Shirley) some 1.
\item True or False: The synagogue in Eshtemoa (As-Samu) was built around the 4 th century. The mosaic floor is decorated with only floral and geometric patterns. The synagogue in Khirbet Susiya (excavated in 1971-72, founded in the end of the 4 th century) has a different answer mosaic panels, the eastern one depicting a Torah shrine, two menorahs, a lulav and an etrog with columns, deer and rams.
\end{enumerate}

\section*{Long Form Response (20 marks)}
\textit{Answer each question with a clear and complete explanation.}
\begin{enumerate}[label=\arabic*.]
\setcounter{enumi}{10}
\item Read the following passage and answer the question: Following the Ulm Campaign, French forces managed to capture Vienna in November. The fall of Vienna provided the French a huge bounty as they captured 100,000 muskets, 500 cannons, and the intact bridges across the Danube. At this critical juncture, both Tsar Alexander I and Holy Roman Emperor Francis II decided to engage Napoleon in battle, despite reservations from some of their subordinates. Napoleon sent his army north in pursuit of the Allies, but then ordered his forces to retreat so that he could feign a grave weakness. Desperate to lure the Allies into battle, Napoleon gave every indication in the days preceding the engagement that the French army was in a pitiful state, even abandoning the dominant Pratzen Heights near the village of Austerlitz. At the Battle of Austerlitz, in Mor... Question: How many muskets did the French capture in the fall of Vienna?
\item Read the following passage and answer the question: In China, urbanization increased as the population grew and as the division of labor grew more complex. Large urban centers, such as Nanjing and Beijing, also contributed to the growth of private industry. In particular, small-scale industries grew up, often specializing in paper, silk, cotton, and porcelain goods. For the most part, however, relatively small urban centers with markets proliferated around the country. Town markets mainly traded food, with some necessary manufactures such as pins or oil. Despite the xenophobia and intellectual introspection characteristic of the increasingly popular new school of neo-Confucianism, China under the early Ming dynasty was not isolated. Foreign trade and other contacts with the outside world, particularly Japan, increased considerably. Chinese ... Question: Name one of the reasons urbanization grew in China.
\end{enumerate}

\vfill
\noindent\textit{End of Paper}
\end{document}